% ============================================================
% 言论新评议（杂志页码 60–63）
% ============================================================
\colchapter{言论新评议}

\section*{让-吕克·戈达尔：新浪潮中的光影人生}

“电影就像面包一样重要，每天都离不开它。我们将电影视作能够获得知识的工具，既是显微镜，又是望远镜。”

——当地时间2022年9月13日，著名导演让-吕克·戈达尔去世，享年91岁。戈达尔是20世纪60年代兴起的法国电影新浪潮运动的标志性人物，从影评人变成导演后，他以《精疲力尽》《狂人皮埃罗》《轻蔑》等作品享誉影坛，也以上述名言激励了全世界范围内无数的后来者，可说是现代电影史上最具影响力的艺术片导演之一。而他的去世标志着法国新浪潮运动的完结，标志着一种电影拍摄理念的终点，也代表了一个时代的落幕。

\section*{自食其力，无关学历}

对高学历人才送外卖，社会应该多一点包容和鼓励。尊重每一份劳动，尊重每一份工作，尊重每一份勇敢尝试的努力。在拼搏的道路上，送外卖是起点不是终点，只要在路上奔走，就不怕到不了未来。

——近日，一则硕士毕业生送外卖的新闻引发网友关注。据报道，中国社科院研究生院2019年硕士毕业生何成曾经从专科生一路“逆袭”读研，却在毕业三年后选择送外卖。何成称，送外卖是自己失业4个月之后的选择，因为没有找到合适的工作才以此过渡。虽然有部分网友质疑研究生送外卖是大材小用，但多数媒体仍对何成的勇气做出如上肯定，并认为“努力的方式不止一种”。

\clearpage
\section*{重庆山火救援：映照家国情怀}

家国情怀，是中国人最朴素、最深沉、最广泛，也是最强大的情感。无论是渡江战役中老百姓划着小舢板冒着枪林弹雨将解放军送过长江，还是汶川地震后15位空降兵在5000米高空的生死“盲跳”，抑或是这次重庆山火中志愿者头灯组成的“星光长城”，在每一个危难关头和紧急时刻，每当国家和人民需要的时候，总有千千万万的人自发地站出来。这就是了不起的中国人，这就是中国的“基本盘”。

——2022年8月中下旬，由于持续的高温干旱，重庆多地突发山林火灾。由于火灾形势严峻，从国家部委到甘肃、云南等兄弟省份，都对重庆进行了大力支援。与此同时，在互联网上，“重庆雄起”的加油打气声此起彼伏，无数令人动容的细节、场景持续刷屏。最终，从第一起山火爆发到各处山火全部扑灭，只用了不到10天的时间，“英雄气”成为社交平台上最火的关键词。而《环球时报》发表文章《英雄气既是重庆的更是中国人的》，以上述评论深情点赞英雄的中国人民。

\subsection*{巧命题}

2020年，美国加州山火持续燃烧70余天；2021年，澳大利亚大火持续210天；2022年，中国重庆山火不到10天就被扑灭。

一场山火，再次证明了中国人民的团结和英雄气概！来自全国各地的消防人员、志愿者和人民子弟兵用血肉铸成火焰烧不透的钢铁长城。有一种凝聚力叫作一方有难八方支援，有一种情怀叫作家国有难舍生向前。这是一个民族该有的样子。

对此，你有什么思考？请写一篇论述类文章加以阐明。

\subsection*{时材立意}

本命题材料通过对比中美澳三国山火的燃烧周期，突出中国人民的英勇无畏以及中华民族强大的凝聚力和浓厚的家国情怀。对此，我们可得出如下立意：

\textbf{①家国情怀是中国人特有的文化基因。}在家国面临危难之际，中国人潜意识里唯一的选择就是尽一切努力去守护。这种家国情怀从何而来？首先，因为共同的生存利益。历史经验告诉我们，国泰才能家齐、民安。其次，中华文化强大的向心力，不断将中国人紧紧凝聚在一起。庞大的文化资产，不断建设着我们的精神世界，使家国思想深入骨髓。最后，数千年的农耕文明让我们对脚下的土地有着深刻的依赖和热爱。这片精耕细作的土地是最宝贵的财产，所以我们要守护和捍卫每一寸土地。

\textbf{②团结一心，展示人民的强大力量。}面对极端高温天气与熊熊燃烧的山火，没有旁观者，都是战斗者。山火的快速扑灭再次展现了中国人民的强大力量，这力量是一点一滴汇聚而来的。如一位摄影记者所说：“那里本没有路，是无数志愿者的脚步；那路本没有灯，是无数志愿者的头灯照亮了重庆儿女用身躯构筑的防火长城。”正是千千万万个普通的救援人员，凝聚成战无不胜的钢铁意志。

\textbf{③天地英雄气，千秋尚凛然。}洪灾、地震、疫情、山火……每当危急关头，这片广袤的土地上都会出现无数挺身而出的平凡英雄。他们选择了勇敢和奉献，就选择了正义和崇高。国泰民安时他们是“凡人”，灾难降临时则立刻变身“英雄”。促使他们成功扮演这样两个角色的原动力只有一个——国家的召唤和人民的期待。

\clearpage
\section*{范文审阅：烈火燔炽灼我衣，与子同袍共偕行}
\auth{文/河北省唐山一中 吴函函（得分：56分）}

这恐怕是“山城”重庆最难熬的一个夏天。2022年8月17日以来，重庆多地接连发生森林火灾。这场山火，侵袭我渝州大地，搅扰我血脉同胞，令人痛心；而平息这场山火，靠的是中华儿女携手共筑的血肉长城，是火一般的爱国热血，令人暖心。

在熊熊燃烧的火焰中，在弥天四起的烟雾中，我们如临深渊、步履维艰，信念却更加坚定，只因看到了你们……

\textbf{偕行之路，有勇气萌生。}我们看到的是驰援重庆的00后消防员薛文——直面火舌，他坦言也会感到紧张，但因为身上的消防服和肩上的责任，他勇敢地坚守着，从未后退。四十多度的高温天气，有无数个如薛文一般的消防员直面火海，用生命保护着生命。是他们与“火魔”持续斗争，才有了普通民众的静好岁月。因为这份同袍之情、这份英勇无畏的民族品质，才有了临渊之时的铿锵步伐。

\textbf{偕行之路，有人间浓浓真情绽放。}我们看到的是无数兄弟省份救援人员，他们与重庆人民携手前行、同此凉热，一道坚不可摧的防线在渝州大地慢慢浮现；我们看到的是主动报名上山清理垃圾，还重庆美丽的千千万万志愿者，他们在用自己的行动诠释“山火无情人人有情”的人间正道。历经几千年的时间沉淀，团结一心、守望相助成为中华民族共同的价值基础，也是中国人最朴素、强大的情感。在每一个危难关头和紧急时刻，这份情义和力量，总是源源不断地涌现。

\textbf{偕行之路，可集思广益。}有外媒曾借重庆山火嘲讽我们灭火方法不先进，却不知我们各部门联动，集结专业力量，灭火方法科学有效，灭火设备先进精准。我们不仅结合风向、地势实施“以火攻火”的传统灭火方法，成功阻挡火势蔓延，还利用“空天地一体化火山监测预警系统”这一“黑科技”指导工作，让行动更为高效。此次扑灭山火，不仅展示了我们专业的力量，更凝聚了我们集体的伟大智慧。

扑灭重庆山火的过程中，我们看到了军民一心迎难而上，还有什么理由不为祖国和民族感到骄傲？我们的土地，是一代一代人披荆斩棘开拓而来的，不光重庆人民爱重庆，全国各族同胞都深爱着这片土地。美哉我少年中国，与天不老，山河无恙！

披子之袍，与子偕行，何愁红日不东升？

\score{得分点：化用《无衣》诗句，鲜明地表达感情，文采斐然。}
\score{得分点：照应标题，明确点出中心论点。}
\score{得分点：素材选择恰当，紧扣分论点中的“勇气”一词。“同袍”二字再扣《无衣》，体现出家国情怀。}
\score{得分点：写灭火的新科技，同时不忘紧扣“偕行”，有任务意识。}
\score{得分点：升华文章主旨，由保卫家园上升到家国情怀，让扑灭重庆山火一事有了更加重大的社会意义。}

\subsection*{时评链接：重庆山火扑救，勇毅逆行筑坚“墙”}

8月26日，经各方救援力量奋力扑救，重庆市北碚区与璧山区交界火场明火全部扑灭。这让连日来为此焦灼、揪心的人们暂时松了一口气。在这场人与火的对峙中，当家园遭遇威胁，同胞置身危难，那些平时默默无闻的人挺身而出，年轻一代奋勇逆行，没有什么得失计算，只觉得义不容辞。为了脚下的土地和心中的信念，他们明知自己作为点点荧火光亮微弱，但仍不放弃闪烁。

还记得那张刷屏的摄影图片吗？燃起大火的森林中，左边是山火燃烧成的红色火线，在不远的右边，一条蓝色光线分外明亮，那是由隔绝带上坚守的志愿者们的蓝色头灯汇聚而成。网友给这张摄影图片起名——“长城”。是的，这座雄伟的长城，我们在1998年长江特大洪水中见识过，在2008年惨烈的汶川地震中见证过，在两年前疫情突袭时甚至都亲自参与构筑过。不论时代如何变迁，社会如何发展，万众一心、众志成城，不畏艰难、守望相助的伟大精神，一直并将继续在中国大地绵延。

（摘自央广网，有删改）

\subsection*{名师态度}

文章标题化用《无衣》，“衣”被大火灼烧后，自然就需要他人与之“同袍”，巧妙地引发读者对英雄的崇敬之情。古诗古人与今文今人结合，感情同频共振，自然流淌。全文采用了“总分总”的结构，“偕行之路”作为筋骨句，分别引领了勇气、真情和智慧三个分论点，每个分论点都与主旨紧密结合，尤其是第三个分论点，很多同学会只强调先进的科技装备与战法，忽略其和主旨“共偕行”的关系。而本文作者却非常成功地用“集思广益”解决了这一问题。

——河北省唐山一中 吴国梁
